\message{ !name(root.tex)}\documentclass[letterpaper,10pt,conference]{ieeeconf}
\overrideIEEEmargins

\usepackage{amsmath,amsfonts,amssymb}
\usepackage{graphicx}
\usepackage{cite}
\usepackage{balance}
\usepackage{algorithm}
\usepackage{algpseudocode}
\usepackage[bookmarks=true,hidelinks]{hyperref}

\hypersetup{
  pdftitle={Arrodes: Recovering Environmental Beliefs from Ergodic Exploration},
  pdfsubject={Inference of an exploring agent's environmental field belief from observed trajectories},
  pdfkeywords={multi-robot exploration, behavior explanation, world-belief inference, ergodic exploration, empirical orthogonal functions, self-organizing maps},
}

\title{\LARGE \bf
Arrodes: Recovering Environmental Beliefs from Ergodic Exploration
}
\author{Anonymous}

\begin{document}

\message{ !name(root.tex) !offset(333) }
\section{Conclusion}

Arrodes extends belief inference to the environmental fields that guide exploratory sampling. SOM states provide a finite set of spatially resolved world hypotheses for exact Bayesian evaluation; EOF coefficients generalize these hypotheses to a continuous space that can represent beliefs outside the SOM's range. The simulated reconstructions show the resulting tradeoff: SOM inference gives an efficient early approximation, whereas EOF inference generally recovers the observed agent's belief more closely once enough of its trajectory has been observed. The relationship between recovery and declining ergodic discrepancy motivates switching between these approaches as evidence accumulates. This ability to reconstruct the model behind an exploration decision is a foundation for coordinating vehicles and explaining their behavior when their environmental knowledge cannot be shared directly.

\balance

\bibliographystyle{IEEEtran}
\bibliography{references}


\message{ !name(root.tex) !offset(330) }

\end{document}
